\chapter{Conclusion}
\label{chap:conclusion}

This thesis set out to carry Free Foil's compile-time scope safety across the LSP boundary into a single, configuration-driven language server, and the result demonstrates that this is achievable.
By parameterizing the server over an opaque \texttt{binder} and bifunctor \texttt{sig} and expressing every operation in terms of \texttt{Bifoldable}, \texttt{Bifunctor}, and a small set of targeted type classes, the library absorbs all tree traversal, scope extension, and name resolution.
A language integrator supplies 15 functions to activate the full feature set, only one of which, the irreducibly language-specific parser invoked by \texttt{buildAsts}, falls outside the \texttt{sig}/\texttt{binder} contract.
Two unrelated languages, LambdaPi and Flan, were served without modifying any core module, confirming that the \texttt{binder}/\texttt{sig} parameterization achieves genuine language agnosticism.

The approach inherits Free Foil's central strength: scope violations become compile-time type errors rather than latent runtime bugs.
Using Haskell's type system to encode the in-scope set as a phantom kind, and using the host language to describe a target language's syntax through a second-order signature, turns a normally bespoke and error-prone layer into reusable, statically checked infrastructure.
The automation and complexity analyses show that this comes at no asymptotic cost: every handler stays within $O(K)$ of the AST size, and go-to-definition and hover, which follow a single root-to-target path, drop to $O(L)$ in the tree depth for fixed-arity grammars.

The framework also has clear boundaries, though they are better understood as what it provides out of the box than as hard limits of the underlying representation.
The current server operates within a single file, models lexical binding, and performs bidirectional type checking rather than constraint-based unification.
None of these are forced by Free Foil itself: because \texttt{S} is only a phantom tag on scopes and the binder is an ordinary type parameter, non-lexical structure can be encoded indirectly, for example as a global namespace as an outermost synthetic binder or as an import that extends the scope with the imported names.
What Free Foil does not supply, unlike scope graphs, is the resolution and dependency machinery itself: visibility, shadowing, qualified lookup, and cross-module name identity must be built by hand, and such encodings recover expressiveness at the cost of some of the by-construction guarantees that motivate the intrinsic approach.
The boundaries are therefore the work Free Foil leaves to the integrator rather than expressive impossibilities, and they delimit the languages the current implementation serves without additional machinery.

Other frameworks that treat name resolution as a first-class, language-independent concern make different trade-offs against the same underlying goal.
Scope graphs and their Statix realization in Spoofax cover non-lexical binding, imports, and rich type systems, but obtain their guarantees through meta-theoretic proof and runtime constraint solving over an extrinsic graph rather than the host type system.
Both formalize a common conception in the world of programming languages, scope, binding, resolution, and types, so that compiler and analyzer behavior can be derived once and reused, and each occupies a distinct point on the axis between expressive breadth and statically checked correctness.
The two families are complementary rather than competing.

A Free-Foil-based language server is therefore already practical for study projects, proof assistants, and simple domain-specific languages, where lexical binding suffices and compiler-checked scope safety is valuable.
Beyond immediate use, it serves as a reference point for how far the intrinsic, type-safe route can be carried into editor tooling, and the limitations identified here, namely cross-file resolution, dependency-driven cache invalidation, incremental type checking, and richer inference, map out concrete directions for further research.
Tools of this kind, which formalize the recurring structures of programming languages and reuse them across implementations, are how the long-term cost of building language tooling is brought down; the work presented here is a step along that path, and a powerful instrument within the scope it addresses.
